\subsection{What an evaluator can change without claiming a new decoder}
\label{sec:p1-recommendations}
The corrected criterion was tested on
already exposed held-person outcomes.
On Shin2017A, source-macro AUC was $0.000307$
lower than with proper predictive NLL selection.
On BNCI2014\_002, it was $0.006819$ higher;
an exploratory interval included zero.
These examples have no independent
positive sign certificate.
An evaluator can nevertheless take
specific, testable precautions:
\begin{enumerate}
\item Label the exact mean-estimation source
in every likelihood score, separating
held-validation-class-centred CovScore from
source-trained-class-mean predictive NLL.
\item Report participant-macro ROC-AUC,
individual harmful effects, and source-selected
parameters beside the density-fitting result;
a change in one is not a substitute for
a change in the others.
\item Do not use held \emph{target}
signals, labels, target-fitted geometry
or performance to choose the source policy
under a claimed zero-target setting.
\item Keep original pre-target protocols
and retrospective density-score diagnostic
analyses in disjoint inferential families.
\item Benchmark prospective selection benefits
against genuine predictive selection and
standard source-macro discrimination, not only
against a historically misnamed covariance score.
\end{enumerate}
The exact decomposed invariant-density term is
a measurement diagnostic for a fixed linear
model, not an externally validated rule
to improve new-person EEG classification
or select a clinical BCI strategy.
